% Rol: Generación de planillas.
\chapter{Generación de planillas}\label{cap:planillas}

\begin{ficharol}{Generación de planillas}{Este capítulo es para usted si atiende a los
  estudiantes y les imprime su planilla de inscripción.}
Con este rol usted puede buscar a los estudiantes cuyo proceso está activo hoy, corregir su
nombre, teléfono, correo y género, y generar o reimprimir su planilla en PDF. No puede cargar
listados ni matricular.
\end{ficharol}

\section{Buscar al estudiante}

Entre a \menu{Admisiones > Generación de Planilla}.

\begin{captura}{planilla-lista}{2732}{1060}
  \anota{1}{46}{151}{738}{192}
  \anota{2}{667}{257}{2065}{552}
  \anota{3}{46}{697}{2686}{1030}
  \anota{4}{2472}{880}{2664}{940}
\end{captura}

\begin{leyenda}
  \item[1] \textbf{Fecha de hoy.} Solo aparecen los estudiantes cuyo proceso está activo hoy.
  \item[2] \textbf{Buscar}: escriba la cédula o parte del nombre y pulse \boton{Buscar}.
        \enlace{Limpiar} vuelve a mostrar a todos.
  \item[3] \textbf{Datos del estudiante.} Si ya tiene planilla, junto a la cédula aparece su
        estado (\emph{pendiente}, \emph{matriculado} o \emph{devuelto}).
  \item[4] \enlace{Seleccionar} abre el detalle del estudiante.
\end{leyenda}

\begin{nota}[¿No aparece el estudiante?]
Si el estudiante no sale en la lista, su proceso no está activo hoy, o su listado todavía no se
ha cargado. Consulte con la persona encargada de la carga de admisiones.
\end{nota}

\section{Revisar y corregir los datos}

Al pulsar \enlace{Seleccionar} se abre el \emph{Detalle de admisión}:

\begin{captura}{planilla-detalle}{2732}{1880}
  \anota{1}{49}{42}{2673}{481}
  \anota{2}{90}{373}{389}{441}
  \anota{3}{90}{723}{1287}{893}
  \anota{4}{1394}{580}{2673}{1819}
  \anota{5}{1871}{1585}{2195}{1685}
\end{captura}

\begin{leyenda}
  \item[1] \textbf{Estudiante}: nombre, teléfono, correo y género. Estos datos sí se pueden
        corregir.
  \item[2] \boton{Guardar datos} guarda las correcciones.
  \item[3] \textbf{Admisión}: estado de la planilla y datos que vienen del listado (cédula,
        carrera, modalidad, etc.). Estos no se pueden cambiar aquí.
  \item[4] \textbf{Campos generados para planilla}: los calcula el sistema.
  \item[5] Botón para \textbf{generar} o \textbf{reimprimir} la planilla.
\end{leyenda}

\begin{pasos}
  \paso{Pregunte al estudiante}{Confirme con el estudiante su nombre, teléfono y correo.}
  \paso{Corrija lo necesario}{Cambie los datos en la sección \emph{Estudiante} y escoja el
        \campo{Genero}.}
  \paso{Guarde}{Pulse \boton{Guardar datos}. Verá el mensaje \emph{«Datos del estudiante
        actualizados.»}}
\end{pasos}

\begin{importante}[La cédula no se cambia aquí]
Si la cédula, la carrera o la modalidad están mal, no genere la planilla: avísele a la persona
de carga de admisiones, porque esos datos vienen del listado oficial.
\end{importante}

\section{Generar la planilla}

\begin{pasos}
  \paso{Revise los datos}{Asegúrese de haber guardado las correcciones.}
  \paso{Genere el PDF}{Pulse \boton{Generar planilla PDF}. La planilla se abre en una pestaña
        nueva del navegador.}
  \paso{Imprima}{Imprima la planilla desde esa pestaña y entréguesela al estudiante.}
\end{pasos}

Así se ve una planilla generada (con datos de prueba):

\pantalla{planilla-pdf-1}[0.8\linewidth]

Al generarla, la admisión queda en estado \textbf{Pendiente}: ya pasa a la cola de
\emph{Revisión de documentos}. Desde ese momento el botón cambia a \boton{Reimprimir PDF}, que
sirve para volver a imprimir la misma planilla.

\begin{consejo}
Si el estudiante pierde la planilla, búsquelo de nuevo, pulse \enlace{Seleccionar} y luego
\boton{Reimprimir PDF}.
\end{consejo}
